\section{Preference Optimization: quién define «una mejor respuesta»}

El prompting puede cambiar una inferencia concreta, pero si se busca que el modelo se incline de forma duradera hacia conductas útiles, seguras o acordes con un estilo determinado, es necesario actualizar los parámetros. La entrada común de la preference optimization son pares chosen/rejected o una reward signal; los métodos difieren en si entrenan un reward model, si muestrean en línea, cómo construyen la baseline y cómo representan las preferencias de una población.

\subsection{RLHF: reward model más policy optimization}

RLHF suele ajustar primero un reward model $r_{\phi}(x,y)$ y después optimiza
\[
\max_{\theta}\;\mathbb{E}_{y\sim\pi_{\theta}(\cdot\mid x)}[r_{\phi}(x,y)]
-\beta D_{\mathrm{KL}}(\pi_{\theta}\|\pi_{\mathrm{ref}}).
\]
El término KL limita cuánto se aleja la policy del reference model. El costo del sistema proviene de la preference collection, el reward model, los rollouts y la estabilidad del RL; el reward hacking muestra que la puntuación sustituta no equivale al valor real para el usuario.


\lecturefigure{slide-078.jpg}{Pipeline de RLHF: comparaciones humanas, reward model y policy optimization}{CS25 V5 official deck, p.~78}

\teachervoice{00:29:13--00:31:32, el ponente separa el prompting del preference training: el primero aporta un scaffold en inference time; el segundo usa human/AI feedback para cambiar los parámetros del modelo. Esta frontera importa más que los acrónimos.}

\subsection{DPO y RLAIF: cambiar el origen de la señal de entrenamiento}

DPO usa directamente preference pairs para optimizar el log-ratio de la policy respecto de la reference:
\[
\mathcal{L}_{\mathrm{DPO}}=-\log\sigma\!\left(\beta\left[
\log\frac{\pi_{\theta}(y^+\mid x)}{\pi_{\mathrm{ref}}(y^+\mid x)}-
\log\frac{\pi_{\theta}(y^-\mid x)}{\pi_{\mathrm{ref}}(y^-\mid x)}
\right]\right).
\]
Prescinde del RL loop explícito con reward model, pero sigue dependiendo de la calidad de los pares, de la reference y de la distribución de preferencias. RLAIF sustituye parte de la anotación humana por un AI judge; reduce costos, pero introduce en la cadena de supervisión el sesgo, la capacidad y la manipulabilidad del judge.


\lecturefigure{slide-079.jpg}{DPO: optimizar la policy directamente a partir de preference pairs}{CS25 V5 official deck, p.~79}


\lecturefigure{slide-080.jpg}{RLAIF: ampliar la anotación de preferencias con AI feedback}{CS25 V5 official deck, p.~80}

\begin{warningbox}{La judge capability es uno de los límites del entrenamiento}
Si el AI judge no puede identificar errores factuales, estrategias de largo plazo o riesgos de seguridad sutiles, escribirá ese ruido en los preference data. Usar un judge más fuerte no resuelve por sí solo el sesgo homogéneo; se requieren human audit, desacuerdo entre varios judges, adversarial prompts y outcome-based checks.
\end{warningbox}

\subsection{GRPO, KTO y heterogeneous preferences}

GRPO construye la advantage a partir del rendimiento relativo de un grupo de sampled responses, con lo que evita parte del costo de un value model independiente. En forma simplificada:
\[
A_i=\frac{r_i-\mu_G}{\sigma_G+\epsilon},
\]
donde $\mu_G,\sigma_G$ son la media y la desviación estándar de la reward dentro del grupo. Depende de la group diversity y de la reward reliability.


\lecturefigure{slide-081.jpg}{GRPO: reward relativa dentro del grupo y actualización de la policy}{CS25 V5 official deck, p.~81}

KTO toma de la prospect theory la loss aversion y admite etiquetas desirable/undesirable en lugar de pairs estrictos; destaca el valor asimétrico de los resultados negativos. Variational Preference Learning, por su parte, modela como variable latente las preferencias de distintos grupos demographic o profiles, y reconoce que «una reward promedio» puede borrar desacuerdos reales.


\lecturefigure{slide-082.jpg}{KTO: loss aversion y señales de preferencia no pareadas}{CS25 V5 official deck, p.~82}


\lecturefigure{slide-083.jpg}{Variational preference learning: personalización y preferencias de poblaciones heterogéneas}{CS25 V5 official deck, p.~83}

\begin{warningbox}{Cinco niveles de validación antes de poner en producción un preference method}
El primer nivel revisa la consistencia de las preference labels, la cobertura de grupos y los atributos sensibles; el segundo mide la held-out pair accuracy, sin tomar la reward-model accuracy como la calidad final del producto; el tercero usa benchmarks multidimensionales para medir por separado helpfulness, truthfulness, harmlessness, style y diversity; el cuarto hace blinded human comparison sobre live-like prompts y registra el disagreement en lugar de promediarlo a la fuerza; el quinto realiza adversarial evaluation para buscar reward hacking, sycophancy, over-refusal y judge manipulation. También hay que comparar el base/reference drift, la KL, la longitud de las respuestas y la calibración. Si un método obtiene un win-rate fuera de línea más alto a costa de colapso de modos, de borrar las preferencias de grupos específicos o de una menor exactitud factual, no representa un avance global. Los loss diagrams de las slides solo definen la trayectoria de optimización; la aceptación debe volver a la distribución real de conductas.
\end{warningbox}

\begin{knowledgebox}{Términos clave: reward, preference y value}
\textbf{Preference label} expresa la elección relativa entre dos respuestas; \textbf{reward model} asigna un escalar a un par entrada-respuesta; \textbf{value/critic} estima el rendimiento futuro; \textbf{reference policy} impide que la actualización se aleje del modelo original; \textbf{judge} puede ser una persona o una IA. Las diferencias entre RLHF, DPO, GRPO y KTO provienen de cómo se combinan estos objetos. Al leer las slides 78--83, conviene dibujar primero el grafo de supervisión y después revisar la loss; memorizar solo las fórmulas hace pasar por alto que la reward distribution, la población que anota y la KL constraint son las que delimitan la conducta.
\end{knowledgebox}

\teachervoice{00:31:32--00:34:13, la clase vincula GRPO, KTO y variational preferences con group-relative reward, loss aversion y demographic heterogeneity, respectivamente. Lo que el docente enfatiza son las objective assumptions, no declarar que algún método sea óptimo en general.}

\begin{table}[H]
\centering
\small
\begin{tabularx}{\textwidth}{L{0.13\textwidth} L{0.20\textwidth} X X}
\toprule
Método & Fuente de supervisión & Componentes que elimina/agrega & Riesgo principal \\
\midrule
RLHF & Preferencias humanas & reward model + RL & reward hacking, rollouts costosos \\
DPO & pair data & Sin RL loop explícito & pair bias, dependencia de la reference \\
RLAIF & AI judge & Mayor capacidad de anotación & judge bias y errores homogéneos \\
GRPO & group rewards & Baseline relativa & Calidad de group/reward \\
KTO & binary desirability & No requiere pairs & Supuestos del value model \\
VPL & profile-conditioned data & Variable latente de preferencia & Privacidad, segmentación y equidad \\
\bottomrule
\end{tabularx}
\caption{Interfaces de supervisión y modos de falla de los métodos de preference optimization.}
\end{table}

\begin{lstlisting}[caption={Esqueleto mínimo de auditoría de un pipeline de preference data}]
prompts = stratified_sample(real_use, safety_edges, long_tail)
responses = generate_candidates(policy, prompts)
labels = collect_preferences(humans, ai_judges, outcomes)
audit_disagreement(labels, by_domain=True, by_group=True)
train_preference_method(labels, reference_model)
validate(helpfulness, factuality, safety, calibration, diversity)
\end{lstlisting}

\subsection{Resumen de la sección}

Lo central de la preference optimization no son los acrónimos, sino quién proporciona la señal, cómo se convierte esa señal en loss y respecto de qué baseline se actualiza el modelo. RLHF, DPO, RLAIF, GRPO, KTO y VPL difieren en costo y en supuestos; ningún método puede sustituir con una sola reward los objetivos reales y multidimensionales del producto.
